\begin{claim}{C.50}{Proved}{Uniform contraction through all the new Lie orbit closures}
In the setting of C.49, let $U_3$ be the union of the open
$G$-orbits having at least three distinct actual lifted
boundary-quotient fiber partitions, and retain $U_2$ from C.37.
Set
\[
                  U_*=U_2\cup U_3,\qquad X_*=\overline{U_*}.
                                                        \tag{C.50.1}
\]
The contractions of C.36 and C.49 extend to a jointly
continuous based contraction through homeomorphisms of $X_*$,
with jointly continuous inverse, commuting with $H$.
This includes arbitrary accumulation of the three-dimensional
homothety orbits.

The uniform algebraic bound is exact. For every continuous
path $k:[0,t]\to\mathcal A_2$ with $k(0)=1$, put
$P=k([0,t])\cup k([0,t])^{-1}$ and $g=k(t)$.
Then
\[
        \{\exp(r\log g):0\leq r\leq1\}
                                      \subseteq P^{1913}.
                                                        \tag{C.50.2}
\]
The identity belongs to $P$, so shorter words can be padded.
For $\R^2$ and $\operatorname{Aff}^+(\R)$ the three-letter
bound of C.37 applies. All word equalities used on an orbit
hold in its specified simply connected source group.
The bound does not identify an additive vector sum with a
noncommutative product: the latter's complete error is
corrected explicitly in the proof.
\end{claim}

\begin{proof}
\emph{Two additive path identities.}
The planar identity of C.37 represents any fraction of the
endpoint of a based path in $\R^2$ by at most three signed
path values. We also require its three-dimensional version:
for a continuous $\gamma:[0,t]\to\R^3$ from $0$ to $v$,
\[
              rv=\sum_{j=1}^{m}\varepsilon_j\gamma(s_j),
 \qquad m\leq5,\quad
       \varepsilon_j\in\{1,-1\},\quad 0\leq s_j\leq t .
                                                        \tag{C.50.3}
\]
For a polygonal path, apply the signed-measure theorem
\cite[Theorem~4.7 and \S4.3.2]{C_GrandoniRaviSinghZenklusen2014}
to its three coordinate derivatives.
The chosen set is a union of at most two circular intervals.
Its integral is a sum of two path differences, with at most
one additional total vector from an interval crossing the
identified endpoint. This gives five signed path values.
Zero terms $\gamma(0)$ permit padding to a nonempty word
when necessary.
The interval count in the halving proof is the same corrected
count used in C.37: two intervals cut at at most three points
give at most five pieces, so one half uses at most two.
For the whole interval, three cuts give at most four pieces,
so one of the two halves uses at most two intervals; these
are also circular intervals.
Complementation, halving, and compactness give every
fraction. Uniform polygonal approximation with the same
endpoints, followed by compactness of the five parameters
and the finite sign set, proves (C.50.3) for arbitrary
continuous paths.

\emph{A maximizing anchor and its exact translation commutators.}
Use the coordinates
\[
        (a,v)(b,w)=(a+b,v+e^a w),\qquad
              v,w\in\R^2,
 \]
and write $k(s)=(a_s,v_s)$.
Let $A=\max_{0\leq s\leq t}|a_s|$.
When $A=0$, the path lies in the translation plane and
the three-letter bound already proves (C.50.2).
For $A>0$, choose $u$ attaining the maximum and let
$h=k(u)$ or $k(u)^{-1}$, choosing its sign so that its
logarithmic scale is $A$.
Thus $h$ is one letter from $P$, not a power of unbounded
length. Conjugate all coordinates by the translation
which sends its unique fixed point to zero.
This is one group automorphism applied to both sides
of every forthcoming identity; it adds no letters to
any word. In the new coordinates $h=(A,0)$ and still
$|a_s|\leq A$. Put
\[
 D=e^A-1,\qquad
 C_s=[h,k(s)]
       =h k(s)h^{-1}k(s)^{-1}
       =(0,Dv_s).
                                                        \tag{C.50.4}
\]
The path $C_s$ is continuous in the translation plane
and begins at the identity.
Every one of its values and inverses uses at most four
letters of $P$.

For any fixed $s$ and $|\beta|\leq B$, where $B$ is a
positive integer, the translation $C_s^\beta$ uses at
most $12B$ letters of $P$.
Indeed, write it as $B$ copies of
$C_s^{\beta/B}$.
For nonnegative $\beta/B\leq1$, the planar identity
applied to the prefix $C|_{[0,s]}$ gives at most three
signed commutator values. For a negative fraction invert
the resulting word. Substituting (C.50.4) costs at most
four letters per commutator.
No continuous choice of the parameters in these words
is required.

\emph{The large-scale case.}
Suppose $A\geq\log2$, so $D\geq1$.
For $0\leq q\leq1$, the translation $T_{qv_s}$ equals
$C_s^{q/D}$ and therefore uses at most twelve letters.
Writing $D_a=(a,0)$ gives
\[
                    D_{a_s}=T_{v_s}^{-1}k(s),
                                                        \tag{C.50.5}
\]
which uses at most thirteen.
Let $a=a_t$ and $v=v_t$. Define
\[
 \theta(a,r)=
 \begin{cases}
  (e^{ra}-1)/(e^a-1),&a\ne0,\\
  r,&a=0.
 \end{cases}
 \qquad
       g^r=(ra,\theta(a,r)v).                           \tag{C.50.6}
\]
The function $\theta$ is continuous and lies in $[0,1]$.
The intermediate value theorem supplies $s$ with
$a_s=ra$, since the path starts at logarithmic scale zero
and ends at $a$.
Then
$g^r=T_{\theta(a,r)v}D_{a_s}$ uses at most
$12+13=25$ letters.

\emph{The small-scale case and its full noncommutative correction.}
Now suppose $0<A\leq\log2$.
Apply (C.50.3) to
$\gamma(s)=(a_s,v_s)$, obtaining
\[
       ra=\sum_{j=1}^{m}\varepsilon_j a_{s_j},\qquad
       rv=\sum_{j=1}^{m}\varepsilon_j v_{s_j},
                         \qquad m\leq5.                 \tag{C.50.7}
\]
Form the actual product
$w=k(s_1)^{\varepsilon_1}\cdots k(s_m)^{\varepsilon_m}$.
For an inverse letter, its translation is
$-e^{-a_{s_j}}v_{s_j}$.
Consequently, with
\[
 p_j=
 \begin{cases}
  \sum_{\ell<j}\varepsilon_\ell a_{s_\ell},
                                       &\varepsilon_j=1,\\
  \sum_{\ell\leq j}\varepsilon_\ell a_{s_\ell},
                                       &\varepsilon_j=-1,
 \end{cases}
                                                        \tag{C.50.8}
\]
the product is exactly
\[
 \begin{split}
 w&=\left(ra,\sum_{j=1}^{m}
                    \varepsilon_j e^{p_j}v_{s_j}\right),\\
 \operatorname{pr}_v(w)
   &=rv+\sum_{j=1}^{m}
             \varepsilon_j(e^{p_j}-1)v_{s_j}.
 \end{split}                                            \tag{C.50.9}
\]
Every $|p_j|$ is at most $mA$.
The pure translation which changes (C.50.9) to (C.50.6)
by left multiplication is
\[
 c=(\theta(a,r)-r)v
       -\sum_{j=1}^{m}
                  \varepsilon_j(e^{p_j}-1)v_{s_j}.
                                                        \tag{C.50.10}
\]
Using the translation commutators (C.50.4), this is the
exact group equation
\[
 g^r=C_t^{\beta_0}
          \left(\prod_{j=1}^{m}C_{s_j}^{\beta_j}\right)w,
 \quad
 \beta_0=\frac{\theta(a,r)-r}{D},\qquad
 \beta_j=-\varepsilon_j\frac{e^{p_j}-1}{D}.
                                                        \tag{C.50.11}
\]
The powers before $w$ commute because they are translations.

Here are uniform bounds, including $a=0$.
The integral formula
\[
        \theta(a,r)=
          \frac{\int_0^r e^{a q}\,dq}
               {\int_0^1 e^{a q}\,dq}                  \tag{C.50.12}
\]
holds at zero as well as away from it.
For $|a|\leq A$, the inequality
$|e^{aq}-1|\leq D$ holds on $0\leq q\leq1$, and the
denominator is at least $e^{-A}\geq1/2$.
Subtracting $r$ in (C.50.12), the numerator is the
difference of an integral over $[0,r]$ and $r$ times
an integral over $[0,1]$, each with integrand
$e^{aq}-1$. Its absolute value is at most $2D$.
Therefore $|\beta_0|\leq4$.
For every other coefficient,
\[
 \frac{|e^{p_j}-1|}{D}
 \leq\frac{e^{mA}-1}{e^A-1}
   =\sum_{\ell=0}^{m-1}e^{\ell A}
   \leq 2^m-1\leq31.                                  \tag{C.50.13}
\]
The same bound covers negative $p_j$ because
$1-e^{-mA}\leq e^{mA}-1$.
The word construction following (C.50.4) now applies with
$B=4$ and $B=31$. Equation (C.50.11) uses at most
\[
                      5+12(4+5\cdot31)=1913             \tag{C.50.14}
\]
original letters. This proves (C.50.2) in all regimes.
The executable certificate accompanying this claim checks
the inverse-letter exponents in (C.50.8), the complete
product correction, the commutator and conjugacy formulas,
and the integer count by exact symbolic arithmetic.
The measure selection and the analytic coefficient bounds
are proved above, independently of that certificate.

\emph{One compact family of actual word isotopies.}
Set $L_0=1913$.
For the original plane loop, all words just proved belong
to the compact parameter family
\[
 W_{\boldsymbol\varepsilon,\boldsymbol s}(u)
       =\prod_{j=1}^{L_0}f_{u s_j}^{\varepsilon_j},
 \qquad
 \varepsilon_j\in\{1,-1\},\quad
 0\leq s_j\leq t\leq1,\quad0\leq u\leq1.
                                                        \tag{C.50.15}
\]
Identity letters are obtained by $s_j=0$.
Continuity of finite multiplication and inversion in the
plane homeomorphism group gives compactness for maps and
inverses, uniformly in $u$.
On a C.49 orbit, the word equality is in its simply
connected actual group $N$.
On a C.36 orbit, use the based path in its specified
universal covering group and the three-letter C.37 bound,
then pad it. Thus the endpoint of the based lift of
(C.50.15) is exactly the required fractional power,
not just a homeomorphism with an unspecified deck shift.

Use $r=1-s$ for the contraction parameter and write
$R^O_{r,t}$ for these fractional actions on an open
orbit in $U_*$. C.36 and C.49 already prove their
individual closure continuity and finite compatibility.
Their groups are metabelian, so
$[[g_1,g_2],[g_3,g_4]]=1$ on $U_*$ for every
$g_i\in G$, and hence on $X_*$.
Every nontrivial nonopen orbit $B\subset X_*$ therefore
has the exact intrinsic image
$M_B=\R$, $\operatorname{Aff}^+(\R)$, or the circle
rotation group, by the path-connected metabelian
argument of C.37.
Define $R^B_{r,t}$ by its unique line-group logarithm
or its based real rotation lift, and use the identity
on point orbits. The based circle lift ends at zero
by C.29. These are the same boundary prescriptions
as in C.36 and C.49.

\emph{Limits, with the bounded residual in the specified cover.}
Take
\[
 x_i\in O_i\subset U_*,\quad x_i\to x\in X_*,
       \qquad (r_i,t_i)\to(r,t).
                                                        \tag{C.50.16}
\]
Choose representations (C.50.15) for
$R^{O_i}_{r_i,t_i}$ and pass to convergent word parameters.
Let $W$ be the limiting plane isotopy, $w=W(1)$,
and $g=f_t$.
If $x$ belongs to an open orbit of $X_*$, then that
orbit already lies in $U_*$ and eventually all $x_i$
lie in it. Its intrinsic continuity applies.
If $x$ is fixed by $G$, every limiting word fixes it.
For a remaining intrinsic line or circle orbit $B=Gx$,
identities on $O_i$ pass to every point $hx\in B$
by evaluating them at $hx_i$ for fixed $h\in G$.
In particular $w|_B$ commutes with $g|_B$.

For a fixed positive integer $n$, set
$m_i=\lfloor nr_i\rfloor$. In each specified simply
connected group, the residual is exactly
\[
       \widetilde w_i^{\,n}\widetilde g_i^{-m_i}
           =\exp\bigl((nr_i-m_i)\log\widetilde g_i\bigr).
                                                        \tag{C.50.17}
\]
Its exponent fraction is in $[0,1)$, so it has another
representation from the same compact word family.
For this fixed $n$, pass to a subsequence on which
$m_i=m_n$ and the residual word isotopies converge
to $Z_n$. Then
\[
       (w|_B)^n(g|_B)^{-m_n}=z_n,\qquad
       |m_n-nr|\leq1,\qquad z_n=Z_n(1)|_B,              \tag{C.50.18}
\]
and all $z_n$ belong to the same compact family of
intrinsic restricted word endpoints.
For a line, logarithms of this compact family are bounded.
In the affine group commuting nonidentity elements belong
to the same exponential one-parameter subgroup; the
identity case gives the same additive logarithm equation.
Thus
\[
          n\log(w|_B)-m_n\log(g|_B)=\log z_n .
 \]
Division by $n$ proves
$\log(w|_B)=r\log(g|_B)$.

For a circle, let $\theta(u)$ be the based real angle
of $f_u|_B$, and put $M=\max_u|\theta(u)|$.
Let $A_W$ be the real endpoint angle of $W|_B$,
$T=\theta(t)$, and $C_n$ that of $Z_n|_B$.
The fixed word length gives $|C_n|\leq L_0M$.
The real, rather than merely circle-valued, equality is
\[
                         nA_W-m_nT=C_n.                \tag{C.50.19}
\]
To prove it, use the plane word isotopies
$W_i(u)^n f_{u t_i}^{-m_i}Z_i(u)^{-1}$.
Their restrictions lift to loops ending at the identity
by (C.50.17). The specified source groups are simply
connected, so these loops are null-homotopic there.
Evaluate at the two points $x_i,hx_i$, where $hx\ne x$.
Their difference-direction winding is zero.
The resulting loops converge uniformly to the analogous
two-point loop at $x,hx$, which stays off the diagonal;
its winding is also zero.
The exact circle comparison of C.29 identifies that
winding, up to the common orientation sign, with
$nA_W-m_nT-C_n$. This proves (C.50.19).
Consequently
\[
                  |A_W-rT|
                       \leq\frac{L_0M+|T|}{n}
                       \quad(n\geq1),
 \]
so $A_W=rT$.
This also covers an identity endpoint on the circle
whose specified real angle $T$ is a nonzero integer.

\emph{Joint continuity and the remaining scope.}
Every convergent subsequence of the compact word
parameters in (C.50.16) has therefore the prescribed
endpoint action at $x$. This proves convergence to
$R^B_{r,t}(x)$ for the entire sequence.
It gives joint continuity from the dense set $U_*$.
For a general sequence in $X_*$, approximate its
$i$th point by a point of $U_*$ with both point distance
and image distance at that fixed parameter less than
$1/i$. The just-proved dense-set limit supplies those
approximations and proves full joint continuity.

Apply the same proof to $f_t^{-1}$.
It has the same generated group, the same open-orbit
families, and the same boundary partitions.
Its fractional actions are the orbitwise inverses,
so the inverse family is jointly continuous as well.
Every member is an orbitwise bijection inverted by
this family and is therefore a homeomorphism of $X_*$.
The equations
$R_{0,t}=1$, $R_{1,t}=f_t|_{X_*}$, and
$R_{r,0}=R_{r,1}=1$ give the based contraction.
The operator transport statements of C.36 and C.49
make $U_*$ and $X_*$ invariant under $H$ and give
commutation on $U_*$. Continuity extends that equality
to $X_*$.

This settles infinite accumulation for the orbit closures
treated by C.49. The remaining global step is a compatible
extension from $X_*$ through the regions with fewer
actual boundary-quotient foliations and their frontiers.
The particular median choice of C.34 cannot be imposed
there without proving compatibility, as C.35 demonstrates.
The general conclusion in C.12 consequently remains open.
\end{proof}
